\section{Real functions and basic examples}

In this course we will often study functions whose inputs and outputs are real numbers. Such functions are called real functions of a real variable. They are central in calculus because they allow us to describe quantities that depend on one real parameter.

A typical notation is
\[
f:D\to\mathbb{R},
\qquad
D\subseteq\mathbb{R}.
\]
This means that the inputs are taken from a set \(D\) of real numbers, and the outputs are real numbers.

\subsection*{Natural domains}

Sometimes the domain is given explicitly. Sometimes it is understood from the formula. When a formula is given and no domain is stated, we often look for the largest set of real inputs for which the formula makes sense. This is called the natural domain.

\begin{examplebox}
Consider
\[
f(x)=2x+1.
\]
This formula makes sense for every real number \(x\). Its natural domain is
\[
\mathbb{R}.
\]
\end{examplebox}

\begin{examplebox}
Consider
\[
g(x)=\frac{1}{x}.
\]
This formula does not make sense for \(x=0\), because division by zero is not defined. Therefore the natural domain is
\[
\mathbb{R}\setminus\{0\}.
\]
\end{examplebox}

\begin{examplebox}
Consider
\[
h(x)=\sqrt{x}.
\]
If we work with real numbers, the expression \(\sqrt{x}\) is real only for \(x\ge0\). Therefore the natural domain is
\[
[0,\infty).
\]
\end{examplebox}

These examples show that the formula is not the whole story. We must also ask where the formula is meaningful.

\subsection*{Linear functions}

A linear function has the form
\[
f(x)=mx+b.
\]
Its graph is a line, provided we view the function as a function of one real variable. The number \(m\) is the slope, and \(b\) is the value at \(x=0\):
\[
f(0)=b.
\]

\begin{examplebox}
Let
\[
f(x)=3x-2.
\]
Then
\[
f(0)=-2,
\qquad
f(1)=1,
\qquad
f(2)=4.
\]
When \(x\) increases by \(1\), the value of the function increases by \(3\). This constant rate of change is the slope.
\end{examplebox}

A linear function is one of the simplest examples of dependence. The output changes at a constant rate as the input changes.

\subsection*{Quadratic functions}

A quadratic function has the form
\[
f(x)=ax^2+bx+c,
\qquad a\ne0.
\]
Its graph is a parabola. The simplest example is
\[
f(x)=x^2.
\]

\begin{examplebox}
For
\[
f(x)=x^2,
\]
we have
\[
f(-2)=4,
\qquad
f(-1)=1,
\qquad
f(0)=0,
\qquad
f(1)=1,
\qquad
f(2)=4.
\]
The values at \(-2\) and \(2\) are the same because squaring removes the sign.
\end{examplebox}

Quadratic functions are already more flexible than linear functions. Their graphs can turn around, and their values need not change at a constant rate.

\subsection*{Reciprocal and square root functions}

Two basic functions show why domains matter:
\[
f(x)=\frac1x
\]
and
\[
g(x)=\sqrt{x}.
\]
The first excludes \(x=0\). The second excludes negative inputs when we work over the real numbers.

\begin{examplebox}
For
\[
f(x)=\frac1x,
\]
we can compute
\[
f(2)=\frac12,
\qquad
f(-4)=-\frac14.
\]
But \(f(0)\) is not defined.
\end{examplebox}

\begin{examplebox}
For
\[
g(x)=\sqrt{x},
\]
we have
\[
g(0)=0,
\qquad
g(4)=2,
\qquad
g(9)=3.
\]
But \(g(-1)\) is not a real number.
\end{examplebox}

These examples are elementary, but they teach an important habit: before evaluating a function, check whether the input belongs to the domain.

\subsection*{Piecewise functions}

A function may use different formulas on different parts of its domain. Such functions are called piecewise-defined functions.

\begin{examplebox}
Define
\[
f(x)=
\begin{cases}
x+1, & x<0,\\
x^2, & x\ge0.
\end{cases}
\]
To compute \(f(-2)\), we use the first rule because \(-2<0\):
\[
f(-2)=-2+1=-1.
\]
To compute \(f(3)\), we use the second rule because \(3\ge0\):
\[
f(3)=3^2=9.
\]
\end{examplebox}

The input decides which rule applies. A piecewise definition is still a function as long as each allowed input receives exactly one output.

\begin{summarybox}
When studying a real function, ask:
\begin{itemize}
\item What is the domain?
\item Is the domain stated explicitly, or should I find the natural domain?
\item Are there inputs excluded by division by zero, square roots, or other restrictions?
\item What type of function am I looking at?
\item If the function is piecewise, which rule applies to the given input?
\end{itemize}
\end{summarybox}

Real functions are the main objects of calculus, but their study begins with reading their domains, formulas, and values carefully.